\usepackage[usenames, dvipsnames]{xcolor}
\usepackage[utf8x]{inputenc}
\usepackage[T1]{fontenc}
\usepackage{amsmath}
\usepackage{tikz}
\usepackage{pgfplots}
% \pgfplotsset{compat=1.9}
\usetikzlibrary{calc,arrows,fadings,decorations.pathreplacing,decorations.markings,patterns,shapes.geometric,circuits.ee.IEC,scopes}
\usepackage{fp}
\usepackage{yhmath}% pour les longueurs d'arc
\usepackage{verbatim}
\usepackage[active,tightpage]{preview}
\PreviewEnvironment{tikzpicture}
\setlength\PreviewBorder{5pt}

%------ mes couleurs -------
\definecolor{fond}{rgb}{1,0.9,0.75}
\definecolor{filet}{rgb}{1,0.6,0}
\definecolor{monOrange}{RGB}{255,157,0}
% \definecolor{monBleu}{RGB}{54,134,193}
\definecolor{monBleu}{rgb}{0.2,0.4,0.6}
\definecolor{monCyan}{RGB}{74,181,247}
\definecolor{monGris}{RGB}{100,100,100}


%------ styles tikz ---------------
\colorlet{darkblue}{blue!50!black} 
\tikzset{>=stealth,inner sep=0pt, outer sep=2pt,}
\tikzset{tiret/.style={gray,dashed}}
\tikzset{doublefleche/.style={|<->|,>=stealth,thin}}
\tikzset{titre/.style={inner sep=0pt, outer sep=0pt,above right,text justified,fill=orange!50}}
\tikzset{bloc/.style={rounded corners=4pt,color=white,ball color=purple,smooth}}
\tikzset{force/.style={->,ultra thick,rounded corners=4pt,color=monBleu,smooth,line cap=round}}
\tikzset{vecteur/.style={->,thick,color=black,smooth}}
\tikzset{verre/.style={draw=SkyBlue,fill=SkyBlue!30}}
\tikzset{axis/.style={thin,gray}}
\tikzset{figure/.style={thick,color=#1,fill=purple, opacity=0.5}}
\tikzset{ressort/.style={very thick,black,smooth}}
\tikzset{eau/.style={draw=black,fill=blue,opacity=0.5}}
\tikzset{trajet/.style={thick,color=blue}}
\tikzset{verre/.style={draw=SkyBlue,fill=SkyBlue!30}}
% style des composants
\tikzset{composant/.style={thick,fill=lightgray,decoration={coil,aspect=0.5,segment length=2mm,amplitude=2mm}}}
\tikzset{bob/.style={decoration={coil,aspect=0.5,segment length=2mm,amplitude=2mm}}}

%-----------------------------------

%-------- patatoide ----------
\newcommand{\patate}[1][fill=white] 
{\draw [#1][preaction={fill=white}] (0,0) .. controls +(0.1,0.2) and +(0.3,0.3) .. (1,0) .. controls +(-0.3,-0.3) and +(-0.05,0.15) .. (1,-1) .. controls +(0.2,-0.6) and +(0.1,-0.2) .. (0,-1) .. controls +(-0.1,0.15) and +(-0.2,-0.4) .. (0,0);}

% --- commande \maSphere{options}{rayon} ---
\newcommand{\maSphere}[2]{
\filldraw[ball color=lightgray!50!white,#1] (0,0) circle(#2);
\draw[dashed,#1] (90:#2)arc(90:-90:{#2/4} and #2);
\draw[dashed,#1] (#2,0)arc(0:180:#2 and {#2/4});
\draw[#1] (#2,0) arc(0:-180:#2 and {#2/4});
\draw[#1] (0,#2) arc(90:270:{#2/4} and #2);
}
% ------ boucle ------
\newcommand{\boucle}[1]{
\draw[#1] (10:3 pt and 10pt) arc(10:350:3pt and 10pt);
\draw[thick,->,#1] (0,0) --++(10pt,0);
}
%------ aimant -------
\newcommand{\aimant}[1]{
\draw[fill, color=black,ultra thick,#1] (0,0.2) rectangle(1,-0.2) node[color=white,midway]{\tiny S};
\draw[fill, color=red,ultra thick,#1] (1,0.2) rectangle(2,-0.2) node[color=white,midway]{\tiny N};}

%------ aimant + champ-------
%------ aimant + champ-------
\newcommand{\aimantPlusChamp}[1]{
\draw[#1,double=monBleu,draw=white](0,3.5)--(0,-3.5);
\draw[<-<,monBleu,#1](0,1)--(0,-1);
\draw[double=monBleu,draw=white,domain=0:{2*pi},samples=50, smooth,#1]plot (xy polar cs:angle=\x r, radius={2*cos(\x r)*cos(\x r)});
\draw[monBleu,<->,domain=.3:{pi-.3},samples=50, smooth,#1]plot (xy polar cs:angle=\x r, radius={2*cos(\x r)*cos(\x r)});
\draw[double=monBleu,draw=white,domain=1.1:{pi-1.1},samples=50, smooth,#1]plot (xy polar cs:angle=\x r, radius={20*cos(\x r)*cos(\x r)});
\draw[monBleu,<->,domain=1.3:{pi-1.3},samples=50, smooth,#1]plot (xy polar cs:angle=\x r, radius={20*cos(\x r)*cos(\x r)});
\draw[double=monBleu,draw=white,domain={pi+1.1}:{2*pi-1.1},samples=50, smooth,#1]plot (xy polar cs:angle=\x r, radius={20*cos(\x r)*cos(\x r)});
\draw[monBleu,>-<,domain={pi+1.3}:{2*pi-1.3},samples=50, smooth,#1]plot (xy polar cs:angle=\x r, radius={20*cos(\x r)*cos(\x r)});
\draw[fill, color=black,ultra thick,#1] (-.15,0) rectangle(.15,-.5) node[color=white,midway]{\tiny S};
\draw[fill, color=red,ultra thick,#1] (-.15,0) rectangle(.15,.5) node[color=white,midway]{\tiny N};
}

% --- support{rayon}{hauteur} ---
\newcommand{\support}[2]{
 \path [top color=black!25,bottom color=white](0,.05*#1/3) ellipse [x radius = #1-.05,y radius = #1/3-.05*#1/3];
 \path [left color=black!25,right color=black!25,middle color= white](-#1,0) -- (-#1,-#2) arc (180:360:#1 and #1/3)  -- (#1,0) arc (360:180:#1 and #1/3);
\foreach \r in {225,315}
	\foreach \i [evaluate = {\s=30}] in {0,2,...,30}
	\fill [black, fill opacity = 1/50](0,0) -- (\r+\s-\i:#1 and #1/3) -- ++(0,-#2) arc (\r+\s-\i:\r-\s+\i:#1 and #1/3) -- ++(0,#2) -- cycle;
\foreach \r in {45,135}
	\foreach \i [evaluate = {\s=30}] in {0,2,...,30}
	\fill [black, fill opacity = 0.02](0,0) -- (\r+\s-\i:#1 and #1/3) arc (\r+\s-\i:\r-\s+\i:#1 and #1/3)  -- cycle;
}

% tire-bouchon \tirebouchon{}
\newcommand{\tirebouchon}[1]{
\draw[black,->,#1](-80:10pt and 3pt) arc(-80:260:10pt and 3pt);
\draw[draw=white,double=black,thick,#1](0,0)--++(0,.5cm);
\draw[thick,draw=black,->,#1](0,.5cm)--++(0,0.1cm);
\draw[decorate,decoration={coil,aspect=0.3,segment length=2pt,amplitude=2pt},#1](0,-1cm+3pt)--++(0,20pt);
\draw[gray,fill=lightgray,#1](-9pt,-3pt-1cm)rectangle(9pt,3pt-1cm);	
}





% commande \resistanceH{option}{nom}  : résistance de longueur unité, centrée en (0,0).
\newcommand{\resistanceH}[2]
{
\draw[#1,fill=white] (-0.5,3pt)--++(1,0)--++(0,-6pt)--++(-1,0)--cycle;
\draw[#1] (0,0) node[above=3pt]{\footnotesize #2};
}

% commande \resistanceV{option}{nom}  : résistance de longueur unité, centrée en (0,0).
\newcommand{\resistanceV}[2]
{
\draw[#1,fill=white] (-3pt,0.5)--++(0,-1)--++(6pt,0)--++(0,1)--cycle;
\draw[#1] (0,0) node[left=3pt]{\footnotesize #2};
}
% \condoH{options}{nom}{valeur} : condensateur horizontal de longueur unité, centrée en (0,0)
\newcommand{\condoH}[3]%
{
\draw[#1,white,very thick] (-3pt,0)--(3pt,0);
\draw[#1,ultra thick] (-3pt,-0.3)--++(0,.6);
\draw[#1,ultra thick] (3pt,-0.3)--++(0,.6);
\draw[#1] (0,-0.3) node[below]{\footnotesize #3};
\draw[#1] (0,0.3) node[above]{\footnotesize#2};
}

% \condoV{options}{nom}{valeur} : condensateur vertical de longueur unité, centrée en (0,0)
\newcommand{\condoV}[3]%
{ 
\draw[#1,white,very thick] (0,-3pt)--(0,3pt);
\draw[#1,ultra thick] (-0.3,-3pt)--++(.6,0);
\draw[#1,ultra thick] (-0.3,3pt)--++(.6,0);
\draw[#1] (0.3,0) node[right]{\footnotesize #3};
\draw[#1] (-0.3,0) node[left]{\footnotesize #2};
}

% \bobineH{options}{nom}{valeur} : bobine horizontal de longueur unité, centrée en (0,0)
\newcommand{\bobineH}[3]%
{
\draw[#1,bob,decorate] (-0.5,0)--++(1,0)node[midway,above=6pt]{\footnotesize #2} node[midway,below=5pt]{\footnotesize #3};
}

% \bobineV{options}{nom}{valeur} : bobine verticale de longueur unité, centrée en (0,0)
\newcommand{\bobineV}[3]%
{
\draw[#1,bob,decorate] (0,-0.5)--++(0,1)node[midway,left=5pt]{\footnotesize #2} node[midway,right=5pt]{\footnotesize #3};
}

% interO{options}{nom} : interrupteurs ouvert horizontale de longueur 0.5 , centré en (0,0)
\newcommand{\interO}[2]%
{
\fill[#1,white] (-10pt,-3pt) rectangle (10pt,3pt);
\draw[#1,-o](-10pt,0)--(-5pt,0);
\draw[#1,-o](10pt,0)--(5pt,0);
\draw[#1,thick] (-7pt,0)++(30:2pt)--++(30:10pt);
\draw[#1] (0,8pt)node{\footnotesize #2};
}

% % interF{options}{nom} : interrupteurs fermé horizontale de longueur 0.5 , centré en (0,0)
\newcommand{\interF}[2]%
{
\fill[#1,white] (-10pt,-3pt) rectangle (10pt,3pt);
\draw[#1,thick] (-5pt,0)--++(10pt,0);
\draw[#1,-o](-10pt,0)--(-5pt,0);
\draw[#1,-o](10pt,0)--(5pt,0);
\draw[#1] (0,8pt)node{\footnotesize #2};
}

% \GBFV{options}{valeur à gauche}{valeur à droite} : GBF verticalde longueur unité, centrée en (0,0)
\newcommand{\GBFV}[3]%
{ 
\draw[#1] (0,-0.5)--++(0,1);
\draw[#1,composant,fill=white] (0,0) circle(0.3);
\draw[#1] (-.2,0) sin (-.1,.1) cos (0,0) sin (.1,-.1) cos (.2,0);
\draw[#1](-3mm,0)node[left]{\footnotesize #2};
\draw[#1](3mm,0)node[right]{\footnotesize #3};
}

% \GBFH{options}{valeur en bas}{valeur en haut} : GBF horizontal de longueur unité, centrée en (0,0)
\newcommand{\GBFH}[3]%
{ 
\draw[#1] (-0.5,0)--++(1,0);
\draw[#1,fill=white] (0,0) circle(0.3);
\draw[#1] (-.2,0) sin (-.1,.1) cos (0,0) sin (.1,-.1) cos (.2,0);
\draw[#1](0,-3mm)node[below]{\footnotesize #2};
\draw[#1](0,3mm)node[above]{\footnotesize #3};
}



% \FEMVL{options}{valeur à gauche}: source de tension verticale de longueur unité, centrée en (0,0)
\newcommand{\FEMVL}[2]%
{ 
\draw[#1,composant,fill=white] (0,0) circle(0.3);
\draw[#1] (0,-0.5)--++(0,1);
\draw[#1,->] (-4mm,-3mm)--(-4mm,3mm)node[midway,left]{\footnotesize #2};
}

% \FEMVR{options}{valeur à droite}: source de tension verticale de longueur unité, centrée en (0,0)
\newcommand{\FEMVR}[2]%
{ 
\draw[#1,composant,fill=white] (0,0) circle(0.3);
\draw[#1] (0,-0.5)--++(0,1);
\draw[#1,->] (4mm,-3mm)--(4mm,3mm)node[midway,right]{\footnotesize #2};
}

% transistor NPN \NPN{options} 
\newcommand{\NPN}[1]
{
\draw[#1,thick,fill=lightgray](0,0) circle(.5);
\draw[#1](-.75,0)--(-0.2,0);
\draw[#1] (-.2,.1)--(60:.5)--++(0,.5);
\draw[#1] (-.2,-.1)--(-60:.5)--++(0,-.5);
\draw[#1,ultra thick, shift={(-.2,0)}](0,-.3)--(0,.3);
\draw[#1,->](-.2,-.1)--(-60:.5);
}

% \transistorH{options} : transistor (commutation) horizontal de longueur unité, centrée en (0,0)
\newcommand{\transistorH}[1]%
{ 
\draw[#1] (-0.5,0)--++(1,0);
\draw[#1,fill=black] (0:2mm)--(120:2mm)--(-120:2mm)--cycle;
\draw[#1](2mm,-2mm)--(2mm,2mm)--++(45:2mm);
\draw[#1](2mm,1mm)--++(45:2mm);
}

% \FEMV{options}{valeur à gauche}: source de tension verticale de longueur unité, centrée en (0,0)
\newcommand{\REVERSEDFEMV}[2]%
{ 
\draw[#1,composant,fill=white] (0,0) circle(0.3);
\draw[#1] (0,-0.5)--++(0,1);
\draw[#1,->] (-4mm,3mm)--(-4mm,-3mm);
\draw[#1](-6mm,0)node{\footnotesize #2};
}

% \CEMV{options}{valeur} : source de courant verticale de longueur unité, centrée en (0,0)
\newcommand{\CEMV}[2]%
{ 
\draw[#1] (0,-0.5)--++(0,1);
\draw[#1,composant,fill=white] (0,0) circle(0.3);
\draw[#1,composant] (-0.3,0)--(0.3,0);
\draw[#1,->] (0,0.3)--++(0,10pt)node[left]{\footnotesize#2};
}

% \CEMH{options}{valeur} : source de courant horizontale de longueur unité, centrée en (0,0)
\newcommand{\CEMH}[2]%
{ 
\draw[#1] (-0.5,0)--++(1,0);
\draw[#1,composant,fill=white] (0,0) circle(0.3);
\draw[#1,composant] (0,-0.3)--(0,0.3);
\draw[#1,->] (0.3,0)--++(10pt,0)node[above]{\footnotesize#2};
}

% \pile{options}{valeur}
\newcommand{\pile}[2]%
{ 
\draw[#1,draw=white,very thick](0,.5)--(0,-.5);
\draw[#1](0,.5)--(0,2pt);
\draw[#1](0,-.5)--(0,-2pt);
\draw[#1] (-10pt,2pt)--++(20pt,0)++(4pt,0)node{\footnotesize +};
\draw[#1,ultra thick] (-3pt,-2pt)--++(6pt,0);
\draw[#1](-10pt,0)node[left]{\footnotesize #2};
}

% \pileH{options}{valeur}
\newcommand{\pileH}[2]%
{ 
\draw[#1,white,very thick](.5,0)--(-.5,0);
\draw[#1](-.5,0)--(-2pt,0);
\draw[#1](.5,0)--(2pt,0);
\draw[#1] (-2pt,-10pt)--++(0,20pt)++(0,4pt)node{\footnotesize +};
\draw[#1,ultra thick] (2pt,-3pt)--++(0,6pt);
\draw[#1](0,-10pt)node[below]{\footnotesize #2};
}

\newcommand{\batterie}[2]%
{ 
\draw[#1,white,very thick](0,.5)--(0,-.5);
\draw[#1](0,.5)--(0,6pt);
\draw[#1](0,-.5)--(0,-6pt);
\draw[#1] (-10pt,6pt)--(10pt,6pt)node[right]{\tiny +};
\draw[#1,ultra thick] (-3pt,2pt)--(3pt,2pt);
\draw[#1](0,2pt)--(0,-2pt);
\draw[#1] (-10pt,-2pt)--(10pt,-2pt);
\draw[#1,ultra thick] (-3pt,-6pt)--(3pt,-6pt);
\draw[#1](-10pt,0)node[left]{\footnotesize #2};
}

% \moteurV{options}{valeur} : moteur vertical de longueur unité centrée
\newcommand{\moteurV}[2]%
{ 
\draw[#1] (0,-0.5)--++(0,1);
\draw[#1,composant,fill=white] (0,0) circle(0.3)node{\small M};
\draw[#1] (-0.3,0)node[left]{\footnotesize #2};
}

% \diodeH{options} : diode horizontale de longueur unité, centrée en (0,0)
\newcommand{\diodeH}[1]%
{ 
\draw[#1] (-0.5,0)--++(1,0);
\draw[#1,composant] (0:2mm)--(120:2mm)--(-120:2mm)--cycle;
\draw[#1,composant](2mm,-2mm)--(2mm,2mm);
}

% \diodeV{options} : diode verticale de longueur unité, centrée en (0,0)
\newcommand{\diodeV}[1]%
{ 
\draw[#1] (0,-0.5)--++(0,1);
\draw[#1,composant] (90:2mm)--(210:2mm)--(-30:2mm)--cycle;
\draw[#1,composant](2mm,2mm)--(-2mm,2mm);
}

% \galva{options} : galvanomètre de longueur unité, centrée en (0,0)
\newcommand{\galva}[1]%
{ 
\draw[thick,#1] (-.5,-.5)rectangle(.5,.5);
% \draw[#1] (0,-.5)node[below]{\footnotesize Galvanomètre};
\draw[#1,fill=lightgray] (0,-.5)++(120:9mm)arc(120:60:9mm)--++(-120:2mm)arc(60:120:7mm)--++(120:2mm);
\draw[#1,thick](0,-.5)++(80:7mm)--++(80:1.5mm);
}

% \voltmetre{options} : voltmètre de longueur unité, centrée en (0,0)
\newcommand{\voltmetre}[1]%
{ 
\draw[thick,fill=white,#1] (-.5,-.25)rectangle(.5,.75);
\draw[#1] (0,.75)node[above]{\footnotesize Voltmetre};
\draw[#1,fill=lightgray] (0,-.25)++(120:9mm)arc(120:60:9mm)--++(-120:2mm)arc(60:120:7mm)--++(120:2mm);
\draw[#1,thick](0,-.25)++(80:7mm)--++(80:1.5mm);
\draw[#1,draw=black,-o](-.5,0)--++(10pt,0);
\draw[#1,draw=black,-o](.5,0)--++(-10pt,0);
\draw[#1](.5cm-10pt,-5pt)--++(4pt,0);
\draw[#1](-.5cm+10pt,-5pt)--++(-4pt,0);
\draw[#1](-.5cm+8pt,-7pt)--++(0,4pt);
}

% --- cylindre{rayon}{hauteur}{couleur1}{couleur2}{option}
\newcommand{\cylindre}[5]{
 \path [#5,top color=#3!25,bottom color=#4](0,.05*#1/3) ellipse [x radius = #1-.05,y radius = #1/3-.05*#1/3];
 \path [#5,left color=#3!25,right color=#3!25,middle color= #4](-#1,0) -- (-#1,-#2) arc (180:360:#1 and #1/3)  -- (#1,0) arc (360:180:#1 and #1/3);
\foreach \i [evaluate = {\s=30}] in {0,2,...,30}{
	\fill [#5,#3, fill opacity = 1/50](0,0) -- (225+\s-\i:#1 and #1/3) -- ++(0,-#2) arc (225+\s-\i:225-\s+\i:#1 and #1/3) -- ++(0,#2) -- cycle;
	\fill [#5,#3, fill opacity = 1/50](0,0) -- (315+\s-\i:#1 and #1/3) -- ++(0,-#2) arc (315+\s-\i:315-\s+\i:#1 and #1/3) -- ++(0,#2) -- cycle;
	}
\foreach \i [evaluate = {\s=30}] in {0,2,...,30}{
	\fill [#5,#3, fill opacity = 0.02](0,0) -- (45+\s-\i:#1 and #1/3) arc (45+\s-\i:45-\s+\i:#1 and #1/3)  -- cycle;
	\fill [#5,#3, fill opacity = 0.02](0,0) -- (135+\s-\i:#1 and #1/3) arc (135+\s-\i:135-\s+\i:#1 and #1/3)  -- cycle;
	}
}





